\documentclass[11pt]{article}
\newcommand{\doctitle}{Positive Definiteness Failures in the
  pmsimstats Covariance Matrix:
  Prevalence, Causes, and Consequences}
\newcommand{\shorttitle}{Positive Definiteness Failures}
\newcommand{\docdate}{2026-04-09}
% pmsimstats-ng standard document preamble
% Usage: % pmsimstats-ng standard document preamble
% Usage: % pmsimstats-ng standard document preamble
% Usage: \input{pmsimstats-preamble} after \documentclass[11pt]{article}
%
% Documents must define before \input:
%   \newcommand{\doctitle}{...}       Full document title
%   \newcommand{\shorttitle}{...}     Short title for header
%   \newcommand{\docdate}{YYYY-MM-DD} Document date

\usepackage[margin=1in]{geometry}
\usepackage{amsmath,amssymb}
\usepackage[T1]{fontenc}
\usepackage{lmodern}
\usepackage{microtype}
\usepackage{graphicx}
\graphicspath{{figures/}}
\usepackage{booktabs}
\usepackage{longtable}
\usepackage{array}
\usepackage{hyperref}
\usepackage{xcolor}
\usepackage{fancyhdr}
\usepackage{parskip}
\usepackage{caption}
\usepackage{enumitem}
\usepackage{verbatim}
\usepackage{listings}

\lstset{
  language=R,
  basicstyle=\small\ttfamily,
  keywordstyle=\color{blue!70!black},
  commentstyle=\color{green!50!black},
  stringstyle=\color{red!60!black},
  breaklines=true,
  frame=single,
  framerule=0.5pt,
  rulecolor=\color{gray!40},
  backgroundcolor=\color{gray!5},
  xleftmargin=1em,
  framexleftmargin=1em,
  columns=flexible
}

\hypersetup{
  colorlinks=true,
  linkcolor=blue!60!black,
  citecolor=green!50!black,
  urlcolor=blue!60!black
}

\pagestyle{fancy}
\fancyhf{}
\lhead{\footnotesize pmsimstats-ng Technical Report}
\rhead{\footnotesize \shorttitle}
\rfoot{\footnotesize \thepage}
\renewcommand{\headrulewidth}{0.4pt}

\title{\doctitle}
\author{pmsimstats team}
\date{\docdate}
 after \documentclass[11pt]{article}
%
% Documents must define before \input:
%   \newcommand{\doctitle}{...}       Full document title
%   \newcommand{\shorttitle}{...}     Short title for header
%   \newcommand{\docdate}{YYYY-MM-DD} Document date

\usepackage[margin=1in]{geometry}
\usepackage{amsmath,amssymb}
\usepackage[T1]{fontenc}
\usepackage{lmodern}
\usepackage{microtype}
\usepackage{graphicx}
\graphicspath{{figures/}}
\usepackage{booktabs}
\usepackage{longtable}
\usepackage{array}
\usepackage{hyperref}
\usepackage{xcolor}
\usepackage{fancyhdr}
\usepackage{parskip}
\usepackage{caption}
\usepackage{enumitem}
\usepackage{verbatim}
\usepackage{listings}

\lstset{
  language=R,
  basicstyle=\small\ttfamily,
  keywordstyle=\color{blue!70!black},
  commentstyle=\color{green!50!black},
  stringstyle=\color{red!60!black},
  breaklines=true,
  frame=single,
  framerule=0.5pt,
  rulecolor=\color{gray!40},
  backgroundcolor=\color{gray!5},
  xleftmargin=1em,
  framexleftmargin=1em,
  columns=flexible
}

\hypersetup{
  colorlinks=true,
  linkcolor=blue!60!black,
  citecolor=green!50!black,
  urlcolor=blue!60!black
}

\pagestyle{fancy}
\fancyhf{}
\lhead{\footnotesize pmsimstats-ng Technical Report}
\rhead{\footnotesize \shorttitle}
\rfoot{\footnotesize \thepage}
\renewcommand{\headrulewidth}{0.4pt}

\title{\doctitle}
\author{pmsimstats team}
\date{\docdate}
 after \documentclass[11pt]{article}
%
% Documents must define before \input:
%   \newcommand{\doctitle}{...}       Full document title
%   \newcommand{\shorttitle}{...}     Short title for header
%   \newcommand{\docdate}{YYYY-MM-DD} Document date

\usepackage[margin=1in]{geometry}
\usepackage{amsmath,amssymb}
\usepackage[T1]{fontenc}
\usepackage{lmodern}
\usepackage{microtype}
\usepackage{graphicx}
\graphicspath{{figures/}}
\usepackage{booktabs}
\usepackage{longtable}
\usepackage{array}
\usepackage{hyperref}
\usepackage{xcolor}
\usepackage{fancyhdr}
\usepackage{parskip}
\usepackage{caption}
\usepackage{enumitem}
\usepackage{verbatim}
\usepackage{listings}

\lstset{
  language=R,
  basicstyle=\small\ttfamily,
  keywordstyle=\color{blue!70!black},
  commentstyle=\color{green!50!black},
  stringstyle=\color{red!60!black},
  breaklines=true,
  frame=single,
  framerule=0.5pt,
  rulecolor=\color{gray!40},
  backgroundcolor=\color{gray!5},
  xleftmargin=1em,
  framexleftmargin=1em,
  columns=flexible
}

\hypersetup{
  colorlinks=true,
  linkcolor=blue!60!black,
  citecolor=green!50!black,
  urlcolor=blue!60!black
}

\pagestyle{fancy}
\fancyhf{}
\lhead{\footnotesize pmsimstats-ng Technical Report}
\rhead{\footnotesize \shorttitle}
\rfoot{\footnotesize \thepage}
\renewcommand{\headrulewidth}{0.4pt}

\title{\doctitle}
\author{pmsimstats team}
\date{\docdate}


\begin{document}
\maketitle
\thispagestyle{fancy}

\begin{abstract}
The \texttt{pmsimstats} data-generating process constructs a
$26 \times 26$ covariance matrix from user-specified correlation
parameters and passes it to \texttt{MASS::mvrnorm()} to draw
simulated participant data. When the resulting matrix is not
positive definite (PD), the code silently projects it to the
nearest PD matrix via
\texttt{corpcor::make.positive.definite()}. The simulation
pipeline was instrumented to measure the frequency and severity
of these failures across four code versions (commits). The results
reveal that PD failure is not a rare edge case but a routine
feature of the simulation: 24.7\% of matrices fail at the
initial commit, rising to 63.0\% after the Ron Thomas edits.
All 162 matrices in the parameter sweep are ill-conditioned
(condition number $> 100$). The failures are concentrated in the
OL+BDC design (83--92\% failure rate) due to abrupt
expectancy-driven variance changes mid-trial. These findings
indicate that the correlation parameter space used in the
publication is not self-consistent and that the PD correction
introduces an uncontrolled perturbation to the intended
covariance structure.
\end{abstract}

\bigskip
\noindent\fbox{\parbox{\dimexpr\textwidth-2\fboxsep-2\fboxrule}{%
\textbf{Architecture scope.} The positive definiteness failures
documented here are specific to Architecture~B (MVN differential
correlation), in which the biomarker-BR correlation parameter
$c_{bm}$ enters the off-diagonal entries of the covariance matrix.
Under Architecture~A (direct mean moderation), the BM-BR
correlation block is absent from $\Sigma$ and the PD constraint on
$c_{bm}$ is eliminated entirely: values up to 0.95 produce valid
covariance matrices across all trial designs. For a comparison of
both DGP architectures, see
\texttt{02-dgp-mean-moderation-vs-mvn.tex}. For implementation
details of Architecture~A, see
\texttt{19-mean-moderation-implementation-notes.tex}.}}
\bigskip

\tableofcontents
\newpage

% =================================================================
\section{Introduction}
% =================================================================

The \texttt{pmsimstats} package simulates clinical trial data by
drawing participant outcomes from a multivariate normal (MVN)
distribution. The MVN requires a positive definite covariance
matrix $\Sigma$. The package constructs $\Sigma$ from a
correlation matrix $R$ and a vector of standard deviations
$\boldsymbol{\sigma}$:
\[
\Sigma = \text{diag}(\boldsymbol{\sigma})
  \cdot R \cdot
  \text{diag}(\boldsymbol{\sigma})
\]

If the resulting $\Sigma$ is not positive definite, the code
applies a correction:

\begin{verbatim}
if (makePositiveDefinite) {
  if (!is.positive.definite(sigma)) {
    sigma <- make.positive.definite(sigma, tol = 1e-3)
  }
}
\end{verbatim}

This correction projects $\Sigma$ to the nearest PD matrix by
setting negative eigenvalues to a small positive tolerance. The
correction is silent: no warning is issued, no diagnostic is
recorded, and the user has no way to know how many matrices
were corrected or how severely the intended covariance structure
was perturbed.

This white paper quantifies the prevalence and severity of PD
failures across the parameter space used in the published
Figure~4, traces their root causes, and evaluates the
consequences for the simulation results.

% =================================================================
\section{Experimental Design}
% =================================================================

\subsection{Instrumentation}

The simulation was instrumented by replacing
\texttt{corpcor::is.positive.definite()} and
\texttt{corpcor::make.positive.definite()} with wrapper
functions that record every call and its result. The wrappers
additionally compute the eigenvalue spectrum and condition
number ($\kappa = \lambda_{\max} / \lambda_{\min}$) for each
matrix checked. No simulation code was modified; the
instrumentation operates at the function-dispatch level.

\subsection{Simulation Parameters}

The parameter sweep is identical to the one used for Figure~4 in
Hendrickson et al.\ (2020):

\begin{itemize}[nosep]
\item \textbf{Trial designs:} OL (1~path), OL+BDC (2~paths),
  CO (2~paths), N-of-1 (4~paths)
\item \textbf{Model parameters:}
  $N \in \{35, 70\}$,
  $c_{bm} \in \{0, 0.3, 0.6\}$,
  $t_{1/2} \in \{0, 0.1, 0.2\}$ weeks
\item \textbf{Fixed correlations:}
  $c_{tv} = c_{pb} = c_{br} = 0.8$,
  $c_{cf1t} = 0.2$,
  $c_{cfct} = 0.1$
\item \textbf{Censoring:} 4 patterns plus no censoring
\end{itemize}

\subsection{Number of PD Checks per Simulation}

The function \texttt{generateSimulatedResults()} calls
\texttt{generateData()} once per path per parameter combination.
Each call constructs one $\Sigma$ matrix and checks it once for
positive definiteness. The total number of checks is
deterministic and independent of the number of replicates
(\texttt{Nreps}), because all $N \times \text{Nreps}$
participants for a given path are drawn in a single
\texttt{mvrnorm()} call:

\begin{table}[h]
\centering
\caption{PD checks per simulation run}
\begin{tabular}{lrr}
\toprule
Design & Paths & $\times$ 18 param sets \\
\midrule
OL & 1 & 18 \\
OL+BDC & 2 & 36 \\
CO & 2 & 36 \\
N-of-1 & 4 & 72 \\
\midrule
\textbf{Total} & & \textbf{162} \\
\bottomrule
\end{tabular}
\end{table}

\subsection{Code Versions Tested}

Four commits spanning the repository history were tested:

\begin{table}[h]
\centering
\caption{Commits tested for PD failure rates}
\begin{tabular}{lll}
\toprule
Commit & Description & Key change \\
\midrule
\texttt{42ac030} & Initial (publication) &
  Step-function BM--BR correlation \\
\texttt{8609f12} & Ron Thomas edits &
  Proportional BM--BR scaling, $s = 2$ \\
\texttt{f6ee86d} & Autocorrelation typo fix &
  Symmetric cross-factor correlation \\
\texttt{f70f86d} & Carryover term added &
  Additional DGP carryover logic \\
\bottomrule
\end{tabular}
\end{table}

% =================================================================
\section{Results}
% =================================================================

\subsection{Overall Failure Rates}

\begin{table}[h]
\centering
\caption{PD failure rates by commit (162 matrices per run)}
\begin{tabular}{llrrr}
\toprule
Commit & Description & Failures & Rate (\%) &
  $\Delta$ from initial \\
\midrule
\texttt{42ac030} & Initial commit &
  40 & 24.7 & -- \\
\texttt{8609f12} & Ron Thomas edits &
  102 & 63.0 & $+38.3$ pp \\
\texttt{f6ee86d} & Autocorr typo fix &
  102 & 63.0 & $+38.3$ pp \\
\texttt{f70f86d} & Carryover added &
  102 & 63.0 & $+38.3$ pp \\
\bottomrule
\end{tabular}
\label{tab:rates}
\end{table}

The Ron Thomas edits (\texttt{8609f12}) more than doubled the
PD failure rate from 24.7\% to 63.0\%. Subsequent commits did
not change the rate, confirming that the additional failures are
caused by the correlation scaling logic introduced at
\texttt{8609f12}, not by the autocorrelation fix or carryover
additions.

\subsection{Condition Number Statistics}

Every matrix in the parameter sweep is ill-conditioned. The
condition number $\kappa = \lambda_{\max} / \lambda_{\min}$
measures numerical stability; values above 100 indicate
ill-conditioning, and values above $10^6$ indicate severe
ill-conditioning.

\begin{table}[h]
\centering
\caption{Condition number statistics across all 162 matrices}
\begin{tabular}{lrrrr}
\toprule
Commit & Median $\kappa$ & $\kappa > 100$ &
  $\kappa > 1{,}000$ & $\kappa > 10^6$ \\
\midrule
\texttt{42ac030} & 229 & 162/162 & 40/162 & 40/162 \\
\texttt{8609f12} & $3.1 \times 10^{18}$ &
  162/162 & 102/162 & 102/162 \\
\texttt{f6ee86d} & $3.1 \times 10^{18}$ &
  162/162 & 102/162 & 102/162 \\
\texttt{f70f86d} & $3.1 \times 10^{18}$ &
  162/162 & 102/162 & 102/162 \\
\bottomrule
\end{tabular}
\end{table}

At the initial commit, 75\% of matrices have $\kappa$ between
100 and 1{,}000 (ill-conditioned but not catastrophic). The
remaining 25\% (the failing matrices) have $\kappa > 10^{18}$,
indicating that the negative eigenvalue is very small relative
to the dominant eigenvalue. After the Ron Thomas edits, the
median $\kappa$ itself is $10^{18}$, reflecting the majority
of matrices now failing PD.

\subsection{Eigenvalue Analysis of Failed Matrices}

All PD failures involve exactly one negative eigenvalue. The
severity ranges from barely negative ($-0.52$) to substantially
negative ($-33.3$):

\begin{table}[h]
\centering
\caption{Minimum eigenvalue in failed matrices}
\begin{tabular}{lrrr}
\toprule
Commit & Worst & Median & Least severe \\
\midrule
\texttt{42ac030} & $-33.26$ & $-15.36$ & $-0.52$ \\
\texttt{8609f12} & $-33.37$ & $-17.11$ & $-0.52$ \\
\texttt{f6ee86d} & $-33.37$ & $-17.11$ & $-0.52$ \\
\texttt{f70f86d} & $-33.37$ & $-17.11$ & $-0.52$ \\
\bottomrule
\end{tabular}
\end{table}

A minimum eigenvalue of $-33$ is not a minor numerical
perturbation. For context, the dominant eigenvalue is
approximately 700 (the sum of the largest variances in the
system). A negative eigenvalue of $-33$ represents approximately
5\% of the dominant eigenvalue, indicating a substantial
structural inconsistency in the specified correlation matrix.

\subsection{Failure Breakdown by Design and Parameter}

The failures are not distributed uniformly across designs. At
the initial commit:

\begin{table}[h]
\centering
\caption{PD failure rate by design and $c_{bm}$ (initial
  commit). Each cell shows failures out of total matrices for
  that design-$c_{bm}$ combination across all $N$ and
  carryover values.}
\begin{tabular}{lrrr}
\toprule
Design & $c_{bm} = 0$ & $c_{bm} = 0.3$ &
  $c_{bm} = 0.6$ \\
\midrule
OL & 0/6 (0\%) & 0/6 (0\%) & 0/6 (0\%) \\
CO & 0/12 (0\%) & 1/12 (8\%) & 1/12 (8\%) \\
N-of-1 & 2/24 (8\%) & 2/24 (8\%) & 2/24 (8\%) \\
OL+BDC & \textbf{10/12 (83\%)} & \textbf{11/12 (92\%)} &
  \textbf{11/12 (92\%)} \\
\bottomrule
\end{tabular}
\label{tab:breakdown}
\end{table}

The OL+BDC design accounts for 32 of the 40 failures (80\%).
The OL design never fails. CO and N-of-1 have low failure rates
that are largely independent of $c_{bm}$, indicating that the
biomarker correlation is not the primary driver of PD failure
for these designs.

% =================================================================
\section{Root Cause Analysis}
% =================================================================

\subsection{Why OL+BDC Fails}

The OL+BDC design has a distinctive feature that none of the
other designs share: the expectancy value changes abruptly from
1.0 (open-label phase, first 4 timepoints) to 0.5 (blinded
discontinuation phase, last 4 timepoints).

The expectancy value directly scales the standard deviation of
the ER (expectancy response) factor:
\[
\sigma_{\text{ER},t} =
  \sigma_{\text{ER}}^{\text{base}} \times e_t
\]
where $e_t$ is the expectancy at timepoint $t$. For OL+BDC:
\[
\sigma_{\text{ER}} =
  [10, 10, 10, 10, 5, 5, 5, 5]
\]

This creates a covariance matrix where the ER block has
variance 100 for the first 4 timepoints and 25 for the last 4,
while the cross-timepoint covariances (driven by the
autocorrelation $c_{pb} = 0.8$) are:
\[
\text{Cov}(\text{ER}_i, \text{ER}_j) =
  0.8 \times \sigma_{\text{ER},i} \times
  \sigma_{\text{ER},j}
\]

When $i$ is in the open-label phase and $j$ is in the blinded
phase:
\[
\text{Cov}(\text{ER}_i, \text{ER}_j) =
  0.8 \times 10 \times 5 = 40
\]

The corresponding correlation in the full covariance matrix is:
\[
\text{Cor}(\text{ER}_i, \text{ER}_j) =
  \frac{40}{\sqrt{100 \times 25}} = 0.8
\]

This is consistent. However, the problem arises from the
interaction between the ER block, the other factor blocks (TV,
BR), and the cross-factor correlations. The compound symmetry
autocorrelation ($\rho = 0.8$) combined with cross-factor
correlations ($c_{cf1t} = 0.2$, $c_{cfct} = 0.1$) and the
heterogeneous ER variances creates a joint covariance structure
that is not guaranteed to be PD.

Intuitively, the abrupt halving of ER variance mid-trial, while
maintaining the same autocorrelation and cross-factor
correlations, forces the correlation matrix into a region where
the implied joint correlations are self-contradictory.

\subsection{Why OL Never Fails}

The OL design has constant expectancy ($e_t = 1$ for all $t$),
producing homogeneous variances within each factor block. The
compound symmetry structure with constant variances is PD by
construction for the autocorrelation and cross-factor
correlation values used (0.8, 0.2, 0.1). The BM--BR
correlations ($c_{bm}$ up to 0.6) do not violate PD in this
homogeneous setting.

\subsection{Why Ron Thomas Edits Increase Failures}

The Ron Thomas edits introduced two changes that affect the
correlation matrix:

\begin{enumerate}
\item \textbf{Timepoint-1 zero correlation.} The BM--BR
  correlation at the first timepoint is always set to 0,
  regardless of treatment status. This creates an additional
  asymmetry in the correlation matrix (the first BR column/row
  has zero correlation with BM while subsequent on-drug
  timepoints have $c_{bm}$).

\item \textbf{Proportional carryover scaling.} Off-drug
  timepoints with carryover residual receive a BM--BR
  correlation proportional to the residual ratio. At the
  carryover half-lives tested (0.1 and 0.2 weeks), these
  correlations are effectively zero (ratios $< 0.001$),
  creating additional near-zero entries in the correlation
  matrix that interact with the existing compound symmetry
  structure.
\end{enumerate}

Both changes increase the heterogeneity of the BM--BR
correlation row/column. At the initial commit, the BM--BR
correlations are either 0 or $c_{bm}$ (a simple binary
pattern). After the edits, they can be 0, a tiny fraction of
$c_{bm}$, or full $c_{bm}$, depending on the timepoint.
This additional heterogeneity pushes more matrices below the
PD boundary.

% =================================================================
\section{Consequences}
% =================================================================

\subsection{What Does \texttt{make.positive.definite()} Do?}

The function \texttt{corpcor::make.positive.definite(m,
tol = 1e-3)} performs the following:

\begin{enumerate}
\item Computes the eigendecomposition: $\Sigma = Q \Lambda Q^T$
\item Replaces any $\lambda_i < \text{tol}$ with $\text{tol}$
\item Reconstructs: $\Sigma^* = Q \Lambda^* Q^T$
\end{enumerate}

This is the ``nearest PD matrix'' in the Frobenius norm sense.
However, it does not preserve the structure of the intended
correlation matrix. Specifically:

\begin{itemize}
\item The diagonal elements of $\Sigma^*$ may differ from the
  intended variances.
\item The off-diagonal elements (covariances) are perturbed
  by an amount proportional to the magnitude of the negative
  eigenvalue.
\item The implied correlation matrix
  $R^* = D^{-1} \Sigma^* D^{-1}$ (where
  $D = \text{diag}(\sqrt{\text{diag}(\Sigma^*)})$) may have
  entries that differ substantially from the intended values.
\end{itemize}

\subsection{Magnitude of the Perturbation}

For the worst-case failures (minimum eigenvalue $\approx -33$),
the correction shifts eigenvalues by up to 33 units. Given that
the total variance in the system (trace of $\Sigma$) is
approximately $\sum_i \sigma_i^2 \approx 3{,}200$, the
perturbation represents roughly 1\% of the total variance.
However, because the perturbation is concentrated in specific
eigenvectors (those associated with the near-zero or negative
eigenvalues), it can disproportionately affect particular
correlation entries (especially the BM--BR correlations and
cross-phase ER correlations that created the non-PD condition
in the first place).

\subsection{Impact on Power Estimates}

The practical impact on power estimates is difficult to quantify
analytically because the perturbation is parameter-dependent and
interacts with the analysis model in complex ways. However,
several qualitative observations can be made:

\begin{enumerate}
\item \textbf{The OL+BDC design is most affected.} With 83--92\%
  of its matrices requiring correction, the OL+BDC power
  estimates are based almost entirely on perturbed covariance
  matrices. The intended correlation structure is not faithfully
  represented.

\item \textbf{The OL design is unaffected.} No corrections are
  needed, so the OL power estimates reflect the intended
  covariance structure exactly.

\item \textbf{Design comparisons are confounded.} When comparing
  power across designs (the paper's primary analysis), the
  OL+BDC results are based on a different effective covariance
  structure than the OL and CO results. This is a systematic
  confound that cannot be resolved by increasing the number
  of replicates.

\item \textbf{The Ron Thomas edits worsen the confound.} By
  increasing the failure rate from 25\% to 63\%, the edits
  ensure that the majority of all simulation runs across all
  designs operate on perturbed covariance matrices.
\end{enumerate}

% =================================================================
\section{Relationship to the Correlation Scaling Artifact}
% =================================================================

This PD failure analysis is related to but distinct from the
correlation scaling artifact described in the companion white
paper. The two issues are:

\begin{enumerate}
\item \textbf{Correlation scaling artifact:} The original code
  assigns full BM--BR correlation to off-drug timepoints with
  negligible carryover residual, creating spurious conditional
  expectations. This is a \emph{modeling choice} issue.

\item \textbf{PD failure:} The specified correlation parameters
  do not produce a valid covariance matrix for certain
  design-parameter combinations, requiring silent correction.
  This is a \emph{parameter validity} issue.
\end{enumerate}

The two issues interact: the correlation scaling rule determines
which BM--BR correlation values are assigned, which in turn
determines whether the resulting matrix is PD. The Ron Thomas
proportional scaling creates more heterogeneous BM--BR
correlation patterns, which increases PD failures. However,
even under the original (simpler) correlation rule, the OL+BDC
design fails PD over 80\% of the time due to the
expectancy-driven variance heterogeneity.

% =================================================================
\section{Recommendations}
% =================================================================

\subsection{Short-Term Mitigations}

\begin{enumerate}
\item \textbf{Log PD failures.} Add a counter to
  \texttt{generateData()} that records how many matrices
  required PD correction, and include this count in the
  simulation output. This allows users to assess the
  prevalence of corrections for any parameter configuration.

\item \textbf{Warn on high failure rates.} If the PD failure
  rate exceeds a threshold (e.g., 10\%), issue a warning
  indicating that the correlation parameters may not be
  self-consistent for the specified design.

\item \textbf{Report the perturbation magnitude.} When a
  correction is applied, compute and log the Frobenius norm of
  the perturbation:
  $\|\Sigma^* - \Sigma\|_F$. This allows users to assess
  whether the correction is negligible or substantial.
\end{enumerate}

\subsection{Medium-Term Solutions}

\begin{enumerate}
\item \textbf{Pre-validate parameter combinations.} Before
  running the simulation, test each unique (design, model
  parameter) combination for PD and reject or flag invalid
  ones. The \texttt{pmsimstats2025} repository implements
  this via \texttt{validate\_parameter\_grid()}.

\item \textbf{Constrain the correlation parameter space.}
  Determine the feasible region of $\{c_{tv}, c_{pb}, c_{br},
  c_{cf1t}, c_{cfct}, c_{bm}\}$ that guarantees PD for each
  design. This can be done numerically by sweeping the
  parameter space and testing PD at each point.

\item \textbf{Address the OL+BDC variance heterogeneity.}
  The OL+BDC design's 83\% failure rate is driven by the
  abrupt expectancy change. Consider whether the compound
  symmetry autocorrelation assumption is appropriate when
  variances change mid-trial, or whether a different
  autocorrelation structure (e.g., damped exponential) would
  be more appropriate.
\end{enumerate}

\subsection{Long-Term Structural Solutions}

\begin{enumerate}
\item \textbf{Cholesky parameterization.} Instead of specifying
  a correlation matrix and checking PD after the fact,
  parameterize the covariance matrix via its Cholesky factor
  $L$ such that $\Sigma = LL^T$ is PD by construction. This
  eliminates the need for PD correction entirely but requires
  reparameterizing the correlation structure.

\item \textbf{Structured correlation models.} Replace the
  ad hoc compound symmetry + BM--BR correlation with a
  structured model (e.g., factor-analytic, spatial
  exponential) that has known PD constraints. This would
  make the feasible parameter region explicit rather than
  requiring numerical validation.

\item \textbf{Separate the correlation parameters from the
  variance parameters.} The current approach specifies
  correlations and variances independently, then combines
  them into a covariance matrix. Because PD depends on the
  joint structure, parameters that are individually reasonable
  can produce an invalid joint specification. A parameterization
  that enforces joint validity would prevent this class of
  errors.
\end{enumerate}

% =================================================================
\section{Conclusions}
% =================================================================

\begin{enumerate}
\item Positive definiteness failure is not a rare edge case in
  the \texttt{pmsimstats} simulation pipeline. At the parameter
  values used in the publication, 24.7\% of covariance matrices
  require silent correction at the initial commit, rising to
  63.0\% after the Ron Thomas edits.

\item All 162 matrices in the parameter sweep are
  ill-conditioned ($\kappa > 100$), and the failed matrices
  have condition numbers exceeding $10^{18}$. The negative
  eigenvalues are substantial (up to $-33$), not mere
  floating-point artifacts.

\item The failures are concentrated in the OL+BDC design
  (83--92\% failure rate) due to the abrupt expectancy-driven
  change in ER variance mid-trial. The OL design never fails.

\item The \texttt{make.positive.definite()} correction silently
  perturbs the intended covariance structure. Because the
  perturbation is design-dependent, cross-design power
  comparisons (the paper's primary analysis) are confounded
  by differences in the effective covariance structure.

\item The Ron Thomas correlation scaling edits increase the
  PD failure rate by introducing additional heterogeneity in
  the BM--BR correlation pattern. This compounds the existing
  OL+BDC failures with new failures across all designs.

\item The correlation parameters used in the publication
  ($c_{tv} = c_{pb} = c_{br} = 0.8$, $c_{cf1t} = 0.2$,
  $c_{cfct} = 0.1$, $c_{bm}$ up to 0.6) are not
  self-consistent for the OL+BDC design. Any simulation
  results for this design should be interpreted with the
  understanding that the covariance structure was
  machine-corrected rather than as specified.
\end{enumerate}

\vfill
\noindent\rule{\textwidth}{0.4pt}
{\footnotesize
Rendered on 2026-04-09 at 18:49 PDT.\\
Source: \textasciitilde/prj/alz/10-pmsimstats-ng/pmsimstats-ng/docs/07-positive-definiteness-failures.tex}

\end{document}
